\section{การวิเคราะห์ระบบ}

จากเป้าหมายหลักคือการเตรียมสื่อสำหรับเผยแพร่โดยลดการเปิดเผยตัวตนของบุคคล ระบบจึงออกแบบให้เริ่มจากการอัปโหลดไฟล์ แสดงผลการตรวจจับในรูปของรายการบุคคลที่จัดกลุ่มแล้ว ให้ผู้ใช้ตัดสินใจว่าจะปกปิดหรือยกเว้นบุคคลใด และแก้ไขบริเวณที่ AI ยังตรวจจับไม่ครอบคลุมก่อนส่งออกไฟล์ ตารางที่ \ref{tab:functional-requirements} แสดงความต้องการเชิงฟังก์ชันของระบบปัจจุบัน

\small
\setlength{\LTcapwidth}{\textwidth}
\begin{longtable}{|p{0.13\textwidth}|p{0.42\textwidth}|p{0.30\textwidth}|}
\caption{ความต้องการเชิงฟังก์ชันของระบบ}
\label{tab:functional-requirements}\\
\hline
\textbf{รหัส} & \textbf{ความต้องการ} & \textbf{สถานะในระบบ} \\
\hline
\endfirsthead
\caption[]{ความต้องการเชิงฟังก์ชันของระบบ (ต่อ)}\\
\hline
\textbf{รหัส} & \textbf{ความต้องการ} & \textbf{สถานะในระบบ} \\
\hline
\endhead
FR-01 & ผู้ใช้อัปโหลดภาพหรือวิดีโอที่รองรับได้ & พัฒนาแล้ว \\
\hline
FR-02 & ระบบแสดงผลตัวอย่างของสื่อที่อัปโหลด & พัฒนาแล้ว \\
\hline
FR-03 & ระบบตรวจจับใบหน้าและบันทึกตำแหน่งตามเฟรม & พัฒนาแล้ว \\
\hline
FR-04 & ระบบรวมใบหน้าที่คาดว่าเป็นบุคคลเดียวกันเป็นกลุ่มบุคคล & พัฒนาแล้ว \\
\hline
FR-05 & ผู้ใช้เลือกเบลอทุกคน เฉพาะคน หรือทุกคนยกเว้นคนที่เลือกได้ & พัฒนาแล้ว \\
\hline
FR-06 & ผู้ใช้กำหนด Gaussian Blur, Pixelation, Black Box, รูปภาพ หรือสติ๊กเกอร์  และปรับความเข้มได้ & พัฒนาแล้ว \\
\hline
FR-07 & ผู้ใช้วาด ย้าย ปรับขนาด ลบ และเลือกติดตามกรอบ Manual Blur ได้ & อยู่ระหว่างพัฒนา \\
\hline
FR-08 & ระบบสร้างตัวอย่างและส่งออกภาพหรือวิดีโอที่ปกปิดข้อมูลแล้วได้ & พัฒนาแล้ว \\
\hline
FR-09 & ระบบส่งไฟล์ผลลัพธ์ไปยัง Cloudinary ได้เมื่อมีข้อมูลรับรอง & พัฒนาแล้ว \\
\hline
FR-10 & ระบบ Livestream ตรวจจับและปกปิดใบหน้าแบบเวลาจริงจากกล้องได้ & พัฒนาแล้ว \\
\hline
FR-11 & ผู้ใช้เลือกตระกูลและขนาดของโมเดลสำหรับ Livestream ได้ & พัฒนาแล้ว \\
\hline
FR-12 & ผู้ใช้เลือกบุคคลได้หลายคน และยกเลิกการเลือกด้วยการกดรายการเดิมซ้ำได้ & พัฒนาแล้ว \\
\hline
FR-13 & ระบบคง Person ID เดิมเมื่อบุคคลออกจากเฟรมชั่วคราวแล้วกลับมา & พัฒนาแล้ว \\
\hline
FR-14 & ระบบแยกภาพต้นฉบับจาก Canvas ที่ใช้แสดงผล เพื่อป้องกันการเบลอซ้ำ & พัฒนาแล้ว \\
\hline
\end{longtable}

ความต้องการที่ไม่ใช่เชิงฟังก์ชัน ได้แก่ การป้องกันพิกัดกรอบไม่ให้เกินขนาดเฟรม การรองรับงานเบื้องหลังระหว่างตรวจจับวิดีโอ การแสดงข้อความแจ้งสถานะการทำงาน และการคงผลการเลือกปกปิดไว้ในฐานข้อมูล ปัจจุบันยังต้องปรับปรุงความปลอดภัยสำหรับการใช้งานจริง เช่น การกำหนด CORS ให้จำกัดโดเมน การบังคับตรวจสอบขนาดไฟล์ที่อัปโหลด และการจัดเก็บสถานะงาน

\section{การออกแบบข้อมูลและส่วนประกอบ}

ฐานข้อมูลของระบบแบ่งข้อมูลเป็นสามกลุ่ม ได้แก่ ข้อมูลสื่อ ผลการวิเคราะห์ใบหน้า และการตั้งค่าปกปิด ตารางที่ \ref{tab:data-model} สรุป entity สำคัญที่กำหนดด้วย SQLAlchemy

% จัดตารางให้ต่อจากหัวข้อโดยคงขนาดตัวอักษรเดิม
\begingroup
\setlength{\intextsep}{0.4\baselineskip}
\renewcommand{\arraystretch}{1.1}
\captionsetup[table]{skip=0.5\baselineskip}
\begin{table}[H]
\centering
\caption{โครงสร้างข้อมูลหลักของระบบ}
\label{tab:data-model}
\begin{tabular}{|p{0.26\textwidth}|p{0.30\textwidth}|p{0.33\textwidth}|}
\hline
\textbf{ตาราง} & \textbf{ข้อมูลสำคัญ} & \textbf{บทบาท} \\
\hline
videos & ชื่อไฟล์, ขนาดเฟรม, FPS, จำนวนเฟรม, สถานะ & เก็บข้อมูลต้นทางของภาพหรือวิดีโอ \\
\hline
preview\_frames & หมายเลขเฟรม เวลา และที่อยู่ของภาพขนาดย่อ & ใช้เลือกและแสดงตัวอย่างเฟรม \\
\hline
detected\_faces & พิกัดกรอบ ค่าความเชื่อมั่น กลุ่มบุคคล และหมายเลขเฟรม & เก็บผลตรวจจับที่ใช้สร้างผลลัพธ์ \\
\hline
identities & ป้ายกำกับ ภาพขนาดย่อ เวกเตอร์ลักษณะเฉลี่ย และจำนวนการปรากฏ & แทนบุคคลแต่ละรายบน Face Map \\
\hline
blur\_selections & รหัสวิดีโอ รหัสกลุ่มบุคคล และสถานะการปกปิด & เก็บการตัดสินใจว่าจะปกปิดบุคคลใด \\
\hline
timeline\_entries & รหัสกลุ่มบุคคล หมายเลขเฟรม และเวลา & แสดงช่วงเวลาที่บุคคลปรากฏ \\
\hline
manual\_blur\_boxes & พิกัด ขนาด เฟรมเริ่มต้น และสถานะการติดตาม & เก็บพื้นที่ปกปิดที่ผู้ใช้กำหนดเอง \\
\hline
\end{tabular}
\end{table}
\endgroup

\begin{samepage}
ส่วนหลัง (backend) แยกความรับผิดชอบเป็นบริการตรวจจับใบหน้า บริการรู้จำใบหน้า บริการติดตามวัตถุ และบริการสร้างผลลัพธ์ การแยกส่วนดังกล่าวช่วยให้เปลี่ยนแบบจำลองหรือปรับวิธีปกปิดได้ แบ่งเป็นคอมโพเนนต์สำหรับอัปโหลดสื่อ เลือกบุคคล แสดง Timeline ควบคุมการปกปิดด้วยตนเอง และส่งออกผลลัพธ์\par
\end{samepage}

\section{การออกแบบส่วนติดต่อผู้ใช้}

ส่วนติดต่อผู้ใช้จัดลำดับการทำงานให้เริ่มจากการอัปโหลดสื่อ เลือกเครื่องมือ ตรวจสอบผลการตรวจจับ แก้ไขพื้นที่ปกปิด และยืนยันผลก่อนส่งออก ภาพ Demo ในส่วนนี้เรียงตามขั้นตอนการใช้งาน

% ใช้ระยะระหว่างกลุ่มภาพที่กระชับ โดยคงช่องว่างภาพ–ชื่อภาพและชื่อภาพ–คำอธิบาย
\begingroup
\setlength{\intextsep}{0.4\baselineskip}

% ==================================================
% TODO IMAGE: ใส่ภาพ Demo หน้าจอหลักและส่วนอัปโหลดไฟล์
% Caption: หน้าจอหลักและส่วนอัปโหลดไฟล์
% ==================================================
\begin{figure}[H]
\centering
\reportscreenshot{\detokenize{images-progress/ภาพ demo หน้าแรก (Home).jpg}}
\caption{หน้าจอหลักและส่วนอัปโหลดไฟล์}
\label{fig:demo-home}
\figuredescription{fig:demo-home}{เป็นหน้าจอเริ่มต้นสำหรับเลือกหรือลากไฟล์ภาพและวิดีโอเข้าสู่พื้นที่อัปโหลดของระบบ}
\end{figure}

% ==================================================
% TODO IMAGE: ใส่ภาพ Demo หน้าจอเลือกเครื่องมือและตั้งค่าการประมวลผล
% Caption: หน้าจอเลือกเครื่องมือและตั้งค่าการประมวลผล
% ==================================================
\begin{figure}[H]
\centering
\reportscreenshot{\detokenize{images-progress/ภาพ demo หน้าการตั้งค่าการประมวลผล(Tools).jpg}}
\caption{หน้าจอเลือกเครื่องมือและตั้งค่าการประมวลผล}
\label{fig:demo-tools}
\figuredescription{fig:demo-tools}{แสดงหน้าจอเลือก AI Auto-Blur, Detect \& Select หรือ Manual Control พร้อมกำหนดรูปแบบและระดับความเข้มของการปกปิด}
\end{figure}

% ==================================================
% TODO IMAGE: ใส่ภาพ Demo หน้าจอการตรวจจับและเลือกบุคคล
% Caption: หน้าจอการตรวจจับและเลือกบุคคลที่ต้องการปกปิด
% ==================================================
\begin{figure}[H]
\centering
\reportscreenshot{\detokenize{images-progress/ภาพ demo หน้าเมนู Detect&Select .jpg}}
\caption{หน้าจอการตรวจจับและเลือกบุคคลที่ต้องการปกปิด}
\label{fig:demo-detect-select}
\figuredescription{fig:demo-detect-select}{แสดงผลการตรวจจับเป็น Face Map เพื่อให้ผู้ใช้ตรวจสอบบุคคลที่พบ เลือกบุคคลที่ต้องการปกปิด และกำหนดรูปแบบการปกปิด}
\end{figure}

% ==================================================
% TODO IMAGE: ใส่ภาพ Demo หน้าจอรายการคีย์เฟรมและ Timeline
% Caption: หน้าจอรายการคีย์เฟรมและช่วงเวลาที่ตรวจพบบุคคล
% ==================================================
\begin{figure}[H]
\centering
\reportscreenshot{\detokenize{images-progress/ภาพ demo หน้าจอแสดงจำนวน Keyframes.jpg}}
\caption{หน้าจอรายการคีย์เฟรมและช่วงเวลาที่ตรวจพบบุคคล}
\label{fig:demo-timeline}
\figuredescription{fig:demo-timeline}{แสดงรายการคีย์เฟรม หมายเลขเฟรม และเวลาที่ตรวจพบบุคคล เพื่อให้ผู้ใช้เลือกตรวจสอบช่วงสำคัญของวิดีโอได้สะดวก}
\end{figure}

% ==================================================
% TODO IMAGE: ใส่ภาพ Demo หน้าจอ Manual Control
% Caption: หน้าจอควบคุมและกำหนดพื้นที่ปกปิดด้วยตนเอง
% ==================================================
\begin{figure}[H]
\centering
\reportscreenshot{\detokenize{images-progress/ภาพ demo เมนู Manual Control.jpg}}
\caption{หน้าจอควบคุมและกำหนดพื้นที่ปกปิดด้วยตนเอง Manual Control}
\label{fig:demo-manual-control}
\figuredescription{fig:demo-manual-control}{แสดงหน้า Manual Control พร้อมพื้นที่ภาพสำหรับวาดกรอบปกปิด ตัวเลือก Track Objects และปุ่ม Save \& Preview สำหรับบันทึกกรอบและตรวจสอบตัวอย่าง}
\end{figure}

% ==================================================
% TODO IMAGE: ใส่ภาพ Demo หน้าจอโหมด AI Auto-Blur
% Caption: หน้าจอการทำงานของโหมด AI Auto-Blur
% ==================================================
\begin{figure}[H]
\centering
\reportscreenshot{\detokenize{images-progress/PreviewAI-AutoBlur.jpg}}
\caption{หน้าจอแสดงความคืบหน้าการประมวลผล AI Auto-Blur}
\label{fig:demo-auto-blur}
\figuredescription{fig:demo-auto-blur}{แสดงหน้าจอ AI Auto-Blur in Progress พร้อมแถบและร้อยละความคืบหน้า ขณะระบบวิเคราะห์เฟรม ตรวจจับใบหน้า และเตรียมผลลัพธ์สำหรับตรวจสอบในขั้นตอนถัดไป}
\end{figure}

% ==================================================
% TODO IMAGE: ใส่ภาพ Demo หน้าจอตรวจสอบผลลัพธ์ก่อนส่งออก
% Caption: หน้าจอตรวจสอบผลลัพธ์หลังประมวลผลก่อนส่งออกไฟล์
% ==================================================
\begin{figure}[H]
\centering
\reportscreenshot{\detokenize{images-progress/ภาพ หน้า Preview Result.jpg}}
\caption{หน้าจอตรวจสอบผลลัพธ์หลังประมวลผลก่อนส่งออกไฟล์}
\label{fig:demo-preview}
\figuredescription{fig:demo-preview}{เป็นหน้าจอตรวจสอบผลหลังประมวลผล ผู้ใช้สามารถเล่นวิดีโอตัวอย่างเพื่อยืนยันว่าพื้นที่เป้าหมายถูกปกปิดครบถ้วนก่อนส่งออกไฟล์}
\end{figure}

\endgroup

\subsection{ส่วนติดต่อผู้ใช้ของ Livestream}

ส่วน Livestream เพิ่มการทดสอบกล้องแบบเวลาจริง โดยผู้ใช้สามารถเลือกตระกูลและขนาดของแบบจำลอง เลือกโหมด \textit{Blur everyone}, \textit{Blur selected person} หรือ \textit{Everyone except selected} และเลือกเทคนิค Gaussian Blur, Pixelate หรือ Black Box พร้อมปรับระดับความเข้มได้ นอกจากนี้ยังรองรับการเลือกบุคคลหลายรายและยกเลิกการเลือกโดยกดปุ่มเดิมซ้ำ

การแสดงผลแยกภาพต้นฉบับออกจากภาพที่เรนเดอร์บน Canvas เพื่อป้องกันการนำภาพที่เบลอแล้วกลับมาเบลอซ้ำ ส่วนการติดตามใช้ Person ID จาก session ร่วมกับ IoU และการกรองตำแหน่ง ทำให้บุคคลที่ออกจากเฟรมชั่วคราวสามารถกลับมาเป็นบุคคลเดิมได้โดยไม่ต้องเลือกใหม่ แผนภาพที่เกี่ยวข้องแสดงในภาพที่ \ref{fig:live-swimlane}, \ref{fig:live-architecture}, \ref{fig:live-activity} และ \ref{fig:live-sequence}

\section{การพัฒนาและการปรับปรุงจากภาคการศึกษาก่อนหน้า}

ตารางที่ \ref{tab:progress-update} แสดงความก้าวหน้าเปรียบเทียบกับสถานะในรายงานเดิม

% ใช้ความกว้างพื้นที่พิมพ์เพื่อลดการตัดบรรทัด โดยคงขนาดฟอนต์
\begingroup
\setlength{\tabcolsep}{4pt}
\renewcommand{\arraystretch}{1.1}
\setlength{\LTpre}{0.4\baselineskip}
\setlength{\LTpost}{0.4\baselineskip}
\captionsetup[table]{skip=0.5\baselineskip}
\small
\setlength{\LTcapwidth}{\textwidth}
\begin{longtable}{|p{0.29\textwidth}|p{0.25\textwidth}|p{0.38\textwidth}|}
\caption{ความก้าวหน้าของระบบในรอบปัจจุบัน}
\label{tab:progress-update}\\
\hline
\textbf{หัวข้อ} & \textbf{สถานะในรายงานเดิม} & \textbf{สถานะปัจจุบัน} \\
\hline
\endfirsthead
\caption[]{ความก้าวหน้าของระบบในรอบปัจจุบัน (ต่อ)}\\
\hline
\textbf{หัวข้อ} & \textbf{สถานะในรายงานเดิม} & \textbf{สถานะปัจจุบัน} \\
\hline
\endhead
ส่วนติดต่อผู้ใช้ & ออกแบบ UI/UX และทดสอบเชื่อม API & มีกระบวนการอัปโหลด Tools, Face Map, Manual Control, Preview และ Export \\
\hline
Selective Blur & พัฒนาแล้ว & จัดกลุ่มบุคคลจากเวกเตอร์ลักษณะ บันทึกตัวเลือกการปกปิด และเลือกปกปิดหรือยกเว้นรายบุคคลได้ \\
\hline
การประมวลผลวิดีโอ & พัฒนาแล้ว & มีกระบวนการตรวจจับแบบเบื้องหลัง การบันทึกตำแหน่งตามเฟรม และการสร้างไฟล์ผลลัพธ์ \\
\hline
Manual Override & อยู่ระหว่างพัฒนา & วาด ย้าย ปรับขนาด ลบ และบันทึกกรอบปกปิด พร้อมตัวเลือกติดตามวัตถุ \\
\hline
รูปแบบการปกปิด & พัฒนาแล้ว & ใช้ Gaussian Blur, Pixelation และ Black Box พร้อมปรับระดับความเข้ม \\
\hline
Livestream & พัฒนาแล้ว & รองรับการเลือกโมเดล โหมดปกปิด การเลือกหลายบุคคล การติดตามด้วย Person ID และการปกปิดบน Canvas แบบเวลาจริง \\
\hline
Preview และ Export & ระบุเป็นหน้าผลลัพธ์ & สร้างตัวอย่างก่อนส่งออกและส่งคืน URL สำหรับดาวน์โหลดหรือเปิดผ่าน Cloudinary \\
\hline
การประเมินเชิงปริมาณ & อยู่ระหว่างพัฒนา & ยังต้องเก็บผลการประเมินประสิทธิภาพและความแม่นยำเพิ่มเติม \\
\hline
\end{longtable}
\endgroup

\subsection{การตรวจจับและสร้าง Face Map}

เมื่อเริ่มกระบวนการตรวจจับ ระบบอ่านวิดีโอทีละเฟรมและเรียกบริการตรวจจับใบหน้า จากนั้นตัดภาพใบหน้าพร้อมขยายขอบเขตบางส่วนเพื่อสร้างเวกเตอร์ลักษณะ ระบบเปรียบเทียบเวกเตอร์ดังกล่าวกับค่าเฉลี่ยของกลุ่มบุคคลที่พบก่อนหน้า หากค่าความคล้ายสูงกว่าค่าเกณฑ์ ระบบจะรวมใบหน้าเข้ากับกลุ่มเดิม หากไม่ใช่จะสร้างกลุ่มบุคคลใหม่ เมื่อเสร็จสิ้น ระบบจะสร้างรายการ Face Map ค่าความเชื่อมั่นเฉลี่ย และ Timeline ของการปรากฏ

สำหรับ pipeline ประมวลผลไฟล์ฝั่งเซิร์ฟเวอร์ การเรียก YOLO เกิดที่คีย์เฟรมที่เลือกหรือเฟรมตามช่วง \texttt{skip\_frames} ส่วนการสร้าง Face Embedding เกิดเฉพาะคีย์เฟรม กรอบที่มีด้านใดด้านหนึ่งเล็กกว่า 32 พิกเซลหรือมี confidence ต่ำกว่า 0.30 จะถูกตัดออกก่อนส่งเข้าขั้นตอนจัดกลุ่ม ดังนั้นการตรวจพบกรอบด้วยโมเดลเพียงอย่างเดียวไม่ได้ยืนยันว่าใบหน้านั้นจะปรากฏใน Face Map

บริการรู้จำใช้ InsightFace/ArcFace และมี DeepFace/ArcFace เป็นทางเลือกสำรองเมื่อโหลดบริการหลักไม่ได้ เวกเตอร์ที่ใช้เปรียบเทียบได้รับการปรับขนาดด้วย L2 normalization แล้วคำนวณ cosine similarity กับตัวอย่างในกลุ่มเดิม การจับคู่ปกติใช้ค่าเกณฑ์ 0.50 และเก็บตัวอย่างเวกเตอร์ของแต่ละกลุ่มได้สูงสุด 8 ตัวอย่าง พร้อมปรับเวกเตอร์เฉลี่ยและมีกระบวนการรวมกลุ่มที่คล้ายกัน ค่าเหล่านี้เป็นค่าที่พบในซอร์สโค้ดปัจจุบัน ยังต้องปรับและยืนยันด้วยชุดทดสอบที่เป็นตัวแทนของการใช้งานจริง

การติดตามจับคู่กรอบด้วย IoU และ Hungarian Algorithm โดย pipeline นี้กำหนด IoU ขั้นต่ำ 0.15 ในเฟรมที่ข้ามการตรวจจับ ตัวติดตามจะประมาณตำแหน่งจากความเร็วเดิม เมื่อได้กรอบตรวจจับใหม่จึงปรับตำแหน่งและรักษารหัสบุคคลของกรอบที่จับคู่ได้ ผลตำแหน่งและรหัสถูกบันทึกเพื่อสร้าง Timeline และใช้ร่วมกับตัวเลือกปกปิดในขั้นตอนเรนเดอร์ การทำงานดังกล่าวเป็นรายละเอียดจากการตรวจโค้ด ส่วนความแม่นยำและความต่อเนื่องยังต้องวัดในการทดลอง

\subsection{การปกปิดแบบเลือกได้และการสร้างผลลัพธ์}

กลุ่มบุคคลที่สร้างขึ้นจะมีสถานะเริ่มต้นเป็นการปกปิด ระบบนำสถานะนี้ไปรวมกับตำแหน่งใบหน้าในทุกเฟรมเมื่อสร้างตัวอย่างหรือส่งออกไฟล์ หากผู้ใช้ต้องการให้ใบหน้าบางรายยังคงมองเห็น สามารถปิดสถานะการปกปิดของบุคคลนั้นได้

\subsection{การปกปิดด้วยตนเองและการติดตาม}

Manual Control แสดงภาพตัวอย่างให้ผู้ใช้ลากกรอบบนวัตถุที่ต้องการปกปิด ตำแหน่งบนหน้าจอจะถูกแปลงกลับเป็นพิกัดจริงของภาพก่อนบันทึก สำหรับวิดีโอ ระบบจะเริ่มตัวติดตามที่เฟรมเริ่มต้นของกรอบ และนำกรอบที่ติดตามได้ไปใช้ปกปิดเฟรมถัดไป

\clearpage

\section{ผลเปรียบเทียบ}

ส่วนนี้แยกผลที่ระบุว่า YOLOv26n (Ours) บนชุดข้อมูล WIDER FACE ออกจากผลอ้างอิง โดยแถว YuNet, RetinaFace และ MTCNN ในตารางที่ \ref{tab:yunet_comparison} นำมาจากงานของ Wu และคณะ \cite{wu2023yunet} ส่วนผลอ้างอิงของแบบจำลองตระกูล YOLO บนชุดข้อมูล COCO แสดงแยกในตารางที่ \ref{tab:yolo-family-reference} เพื่อไม่ให้เปรียบเทียบตัวชี้วัดจากคนละชุดข้อมูลโดยตรง

\begin{table}[H]
\centering
\small
\caption[ผลการเปรียบเทียบโมเดลตรวจจับใบหน้าและผลอ้างอิงแบบจำลอง YOLO ทุกขนาด]{ผลการเปรียบเทียบโมเดลตรวจจับใบหน้าและผลอ้างอิงแบบจำลอง YOLO ทุกขนาด~\cite{wu2023yunet}}
\label{tab:yunet_comparison}
\resizebox{\textwidth}{!}{%
\begin{tabular}{c l r r c c c c}
\toprule
\textbf{Input size} &
\textbf{Methods} &
\textbf{\#FLOPs (M)} &
\textbf{$AP_{easy}$} &
\textbf{$AP_{medium}$} &
\textbf{$AP_{hard}$} &
\textbf{Latency (ms)}
\\
\midrule

% ==================== 320 x 320 ====================
\multirow{4}{*}{320 $\times$ 320}

& YOLOv26n (Ours)
& 194
& 0.894
& 0.8677
& 0.7257
& 7.2
\\

& YuNet
& 149
& 0.836
& 0.747
& 0.395
& 2.2
\\

& RetinaFace
& 11\,070
& 0.868
& 0.742
& 0.341
& 49.1
\\

& MTCNN
& 3\,359
& 0.923
& 0.862
& 0.504
& 17.3
\\

\midrule

% ==================== 640 x 640 ====================
\multirow{4}{*}{640 $\times$ 640}

& YOLOv26n (Ours)
& 692
& 0.907
& 0.894
& 0.880
& 17.8
\\

& YuNet
& 595
& 0.899
& 0.869
& 0.691
& 20.4
\\

& RetinaFace
& 44\,260
& 0.943
& 0.908
& 0.659
& 232.7
\\

& MTCNN
& 11\,545
& 0.949
& 0.935
& 0.814
& 104.0
\\

\midrule
\bottomrule
\end{tabular}%
}

\end{table}

ตาราง \ref{tab:yunet_comparison} แสดงว่าแบบจำลองแต่ละชนิดมีข้อแลกเปลี่ยนต่างกัน ที่ขนาด 320 พิกเซล YuNet ใช้เวลาน้อยกว่า YOLOv26n ขณะที่ YOLOv26n มีค่า AP ในกลุ่ม Medium และ Hard สูงกว่า การเลือกแบบจำลองจึงต้องพิจารณาทั้งความละเอียด ความแม่นยำ และเวลา โดยไม่สรุปว่าแบบจำลองหนึ่งดีกว่าทุกเงื่อนไข ทั้งนี้ ควรแนบรายละเอียดชุดทดสอบ น้ำหนักแบบจำลอง และฮาร์ดแวร์ของผลที่ระบุว่า Ours ก่อนใช้ยืนยันประสิทธิภาพของระบบ

\clearpage
\small
\setlength{\LTcapwidth}{\textwidth}
\begin{longtable}{c c c r r r r}
\caption[ผลแบบจำลอง YOLO บน COCO ที่ความละเอียด 640 พิกเซล]{ผลแบบจำลอง YOLO บน COCO ที่ความละเอียด 640 พิกเซล~\cite{ultralytics2023yolov8,ultralytics2024yolo11,ultralytics2026yolo26docs}}
\label{tab:yolo-family-reference}\\
\toprule
\textbf{Input size} & \textbf{Family} & \textbf{Size} & \textbf{mAP$_{50:95}$} & \textbf{CPU ONNX (ms)} & \textbf{Params (M)} & \textbf{FLOPs (B)} \\
\midrule
\multirow{15}{*}{640 $\times$ 640} & \multirow{5}{*}{YOLOv8} & n & 37.3 & 80.4 & 3.2 & 8.7 \\
 & & s & 44.9 & 128.4 & 11.2 & 28.6 \\
 & & m & 50.2 & 234.7 & 25.9 & 78.9 \\
 & & l & 52.9 & 375.2 & 43.7 & 165.1 \\
 & & x & 53.9 & 479.1 & 68.2 & 257.8 \\
\midrule
& \multirow{5}{*}{YOLO11} & n & 39.5 & 56.1 & 2.6 & 6.5 \\
 & & s & 47.0 & 90.0 & 9.4 & 21.6 \\
 & & m & 51.5 & 183.2 & 20.1 & 68.1 \\
 & & l & 53.4 & 238.6 & 25.3 & 87.2 \\
 & & x & 54.7 & 462.8 & 56.9 & 195.3 \\
\midrule
& \multirow{5}{*}{YOLO26} & n & 40.9 & 38.9 & 2.4 & 5.5 \\
 & & s & 48.6 & 87.2 & 9.5 & 20.9 \\
 & & m & 53.1 & 220.0 & 20.4 & 68.4 \\
 & & l & 55.0 & 286.2 & 24.8 & 86.8 \\
 & & x & 57.5 & 525.8 & 55.7 & 194.4 \\
\bottomrule
\end{longtable}



ผลลัพธ์ตามตาราง \ref{tab:yolo-family-reference} เป็นผลอ้างอิงจากเอกสารของ Ultralytics \cite{ultralytics2023yolov8,ultralytics2024yolo11,ultralytics2026yolo26docs} ใช้เพื่อแสดงแนวโน้มของแต่ละขนาดเท่านั้น ค่า COCO mAP ไม่สามารถนำไปแทนค่า AP ของงาน Face Detection โดยตรง
